\section{My impact}

% ============================================================
% 3a. CONTRIBUTIONS DRIVING THE TEAM FORWARD
% ============================================================

\subsection{3a. How my contributions drove the team forward}

The question of impact is, in the rubric's language, whether my contributions \textit{drove the team forward} --- and in mine, whether the things I built unblocked other people's work or only my own.

The honest thesis of this section is the one I wrote into our group complaint letter: \textit{the simulation framework we built ourselves is the only reason this project produced any verifiable engineering output at all.} I want to defend that claim with mechanism, not assertion, because it is exactly the kind of line that at Merit-level reads as self-flattery and at Distinction-level requires evidence.

\textbf{The Kolb cycle that defined this project.} I want to walk through a full Kolb (1984) cycle here, because I used it deliberately rather than retrospectively, and it is the single piece of experiential learning that most changed my engineering practice.

\textit{Concrete experience}: Three collapsed flight days. The 2026-03-11 weather cancellation, where I stood on a wet field in North Somerset with a drone that would not fly. The 2026-03-30 hardware session that produced zero airborne minutes because the calibration data was corrupt. And a third day that ended with our pilot standing on a wet field while Edward and I debugged \texttt{calibration\_data.npz} --- an empty file that caused the remap matrices to mangle every frame, making the AI see noise instead of images. Three scheduled flight days; zero seconds of real flight data; approximately eight weeks of calendar time consumed.

\textit{Reflective observation}: After the first cancellation, the observation was harder than ``we lost a day.'' The team's entire definition of ``progress'' had silently assumed hardware access we did not have. Robin could not test his flight state machine. Edward could not validate his Cube interface. Demetro could not generate real-world detection data for his presentation. The only person who was not fully blocked was me, because I had a laptop with SITL. That asymmetry --- I could work, they couldn't --- was the data point that mattered.

\textit{Abstract conceptualisation}: If \texttt{vision.py} exposed a single interface (\texttt{detect\_in\_image(frame)} returning \texttt{(found, x, y, conf)}) and auto-selected between Ultralytics on laptop and TFLite on Pi, and if \texttt{config.py} auto-detected its hardware, then the same state machine, geofence, and search pattern could run in both places. Anything that ran in simulation would, in principle, run on hardware. The key design decision --- and I want to name it as a decision, not an accident --- was to make \texttt{vision.py} the ONLY file that touches CV/AI. Everything else calls the same function. You can change the model, the preprocessing, the backend, without touching any other file. That modularity was not an aesthetic choice; it was a risk mitigation. In Graham's (2026) risk management framework, it reduced both the likelihood of a platform-incompatibility failure (by abstracting the backend) and the consequence (by isolating the blast radius to one file).

\textit{Active experimentation}: Six weeks of hammering on that premise. \texttt{simple\_simulator.py} grew to 2508 lines with simulated GPS drift (configurable via \texttt{--gps-drift}), camera shake (\texttt{--shake}), CV throttle, and a TFLite backend path. \texttt{main.py} was refactored from 1411 to 754 lines across six extracted modules with 146 pytest tests. The progressive test ladder (\texttt{tests/flight/0a}--\texttt{4}) plus 127 unit tests gave me a way to know, before any field day, which scripts could survive contact with hardware. When we finally ran on the Pi on 2026-03-31, the delta between simulation and hardware was thin: a bbox-coordinate bug where TFLite pixel coords (0--640) were treated as normalised values (0--1), causing green detection boxes to render off-screen (commit \texttt{3d62e6a}); a double-scaling fix where \texttt{detect\_in\_image} returned pixel coords but \texttt{passive\_watch.py} multiplied by width and height again, offsetting the GPS estimate by more than 100 metres (commit \texttt{7eb7cc8}); and a deleted calibration file that had been silently corrupting frames. The system came up.

The Kolb cycle taught me something I had not previously articulated: simulation-first is not a personal preference. It is a claim about what the engineering discipline owes to a project that lacks hardware access. And the claim has limits. The simulator could not predict motion blur at 5\,m/s with a real IMX296 sensor. It could not find the descent stall bug we eventually accepted as a SITL artefact. It could not tell us what GPS timing lag (100--200\,ms latency, causing $\sim$1\,m error at 5\,m/s) would do to a multi-pass target lock --- that insight came from DJI video analysis, after the fact. The statement my simulation made is true, and the limits of that statement are also true. Sim-to-real parity is an argument I can only win about the things I modelled, and I was disciplined about that only because I had been forced to be.

\textbf{Three other contributions, compactly.}

\textit{Unblocking Robin.} I built Robin a dedicated HTTP integration path (\texttt{passive\_watch\_2.py}, commit \texttt{a5b1ab7}) with a JSON schema, \texttt{cv\_mode} polling, and a \texttt{--smart-clear} flag so his mission script could poll vision results without touching the CV code. I wrote a README specifically for him: exact CLI command, JSON schema, parsing examples (commit \texttt{885cee4}). This was anticipating a teammate's needs instead of dumping code. The \texttt{docs/VISION\_PIPELINE\_FOR\_COLLEAGUE.md} and \texttt{docs/COLLEAGUE\_CHECKLIST.md} --- 43 items from cloning the repo to running the first detection --- were written for whoever came after Robin.

\textit{Unblocking Demetro.} I built Demetro's two FDR slides from scratch (commit \texttt{32fb0c5}) and wrote full speaking scripts: 246 + 219 = 465 words, 3.1 minutes total (commits \texttt{b2a9cd4}, \texttt{ba667f2}). I iterated 4--5 times, including a teleprompter sidebar. This was team-first behaviour: I wrote and polished my teammate's presentation materials, not my own.

\textit{Driving external communication.} I led the team's external communication about hardware failures via the six-iteration complaint letter (\texttt{docs/COMPLAINT\_LETTER.md}). The letter was signed by all five team members. The tone was deliberately restrained: ``We are not assigning personal blame. We are documenting what happened.'' Commits \texttt{26e3255} through \texttt{4b0fd91} show six iterations on the tone, phrasing, and evidence. This was a decision to channel team frustration into a professional artefact rather than letting it dissipate as corridor complaints.

\textbf{Evaluating my own effectiveness (M16.d).} Against the criterion of ``did my contributions unblock teammates or only advance my own work,'' the judgement is mixed. The simulation framework was genuinely team-enabling: it let Robin develop his flight script, Edward validate his Cube interface, and Demetro generate synthetic data, all without a drone. The handoff documentation worked: Robin integrated in wk20 without clarifying questions. But the simulator itself (\texttt{simple\_simulator.py}, 2508 lines) was edited by zero teammates across the entire project. The unit tests I wrote (127 functions across 6 files) were run by zero teammates. The main.py refactor (1411 to 754 lines) was reviewed by zero teammates. The cause is not that these artefacts were bad --- it is that I built them as solo infrastructure rather than collaborative tools. The lesson: I was effective at producing artefacts but ineffective at producing shared ownership, and those are different kinds of effectiveness.

\textbf{Evaluating my team's effectiveness (M16.d-team).} Against the criterion of ``did the team meet its obligations and goals,'' the judgement is: partially. We met 12/12 requirements in simulation but only 4/12 on real hardware by the final flight day. The gap was caused by three factors: hardware unavailability (external, uncontrollable), coordination failures (internal, our responsibility), and insufficient investment in team processes (internal, largely my responsibility as the person who set the technical pace). Our Tuckman progression --- brief Forming, extended Storming, late Norming, questionable Performing --- reflects a team that invested too little time in managing itself, exactly the pattern Graham warns about: ``teams that don't bother to think about how they work will have to spend even more time trying to manage the team.'' The cause: we treated process as overhead rather than infrastructure, and the person who should have advocated most loudly for process --- me, the de facto technical lead --- was the person most allergic to it.

\textbf{Six communication methods evaluated (M17.d).} I used six communication methods across the project, each with a different audience:

\begin{enumerate}[nosep]
\item \textbf{Markdown handoff docs} (Robin, wk19) --- 3000-word technical pipeline documentation
\item \textbf{In-person pair walkthrough} (Edward, wk20) --- sitting together debugging the Cube serial connection
\item \textbf{Rubric-stripped plain-language bug report} (PM, wk21) --- ``the drone is waiting for a message that never comes'' instead of ``ACK race in \texttt{MAV\_CMD\_NAV\_TAKEOFF}''
\item \textbf{Minimalist button UI} (\textbf{the pilot}, wk22) --- four buttons (Y / N / M / L) mapped to the RC thumb-reach pattern
\item \textbf{Five-draft FDR speaking script} (Demetro, wk16) --- iterating until the words matched his speaking cadence
\item \textbf{Drafted complaint letter} (\textbf{Steve, the course organiser} --- a non-engineering audience) --- formal, factual, signed by all five team members
\end{enumerate}

Against a single criterion --- whether the recipient acted within 48 hours without coming back with clarifying questions --- methods 1 and 4 hit the threshold directly. Method 2 hit it because pair walkthrough compresses the whole clarification loop into the session itself. Method 3 hit it because I stripped two layers of jargon (``ACK race in \texttt{MAV\_CMD\_NAV\_TAKEOFF}'' became ``the drone is waiting for a message that never comes''). Method 5 hit it after five drafts because my first three were in my voice, not Demetro's. Method 6 hit it because the complaint letter had to survive a reader with no engineering background.

The lesson I extracted: communication effectiveness is not about clarity --- it is about \textit{cost asymmetry}. I spent roughly twenty hours on the Markdown doc (method 1) and Robin spent four hours integrating. Pair walkthroughs (method 2) inverted that ratio. The same author-hour budget produces dramatically different behaviour change depending on which method the hours go into. One reading is that I should have done more pair work. Another is that the documentation was a time-shifted investment that paid off weeks later when Robin needed it. I conclude that both are true, and that the right strategy is to pair-program for urgent unblocking and document for future-proofing --- but I should have been explicit about which mode I was in, rather than defaulting to documentation because it is my comfort zone.

The non-technical audience adaptation (method 6, the complaint letter to Steve) deserves specific attention. The audience was the course organiser --- an academic, not an engineer --- and the content was a factual account of hardware failures and their impact on learning outcomes. The adaptations I made: (1) removed all code references and commit hashes from the body text; (2) framed the argument around learning objectives rather than technical capabilities (``we could not evidence R01--R12 on real hardware'' rather than ``the geofence never ran on the Cube''); (3) structured the letter with a clear problem-impact-request pattern rather than a chronological narrative; and (4) showed a draft to Robin, who caught three sentences that were ``technically accurate but personally sharp'' and suggested softer alternatives. The letter was received and acknowledged, and the course organiser adjusted the assessment criteria for teams affected by hardware failures.

% ============================================================
% 3b. ADD OR RE-FOCUS IN FUTURE
% ============================================================

\subsection{3b. What I would add or re-focus on in future projects}

Against the twelve project requirements R01--R12, our system hit 12/12 in simulation but only 4/12 on real hardware by flight day. Three of the eight gaps were bench-test items rather than engineering failures --- evidence, not apology, that coordination was our limiting factor, not effort. But I want to be precise about the nature of the gap, because ``coordination'' is the kind of word that sounds like an explanation and is actually a label.

The specific gap: I mistook technical leadership for team leadership. I now see they are orthogonal --- you can build great tools that do not enable anyone --- and in this project the gap showed up as a metric: the number of times a teammate other than me edited \texttt{simple\_simulator.py} or \texttt{main.py} was zero. That is not a coordination failure; it is an ownership failure. I owned too much, shared too little, and called it diligence.

Graham's (2026) risk management lecture introduces three pillars of obligation: moral, economic, and legal. I want to apply a parallel framework to my impact assessment. Morally, I have an obligation to build things that help the team, not just things that help me look productive. Economically, the team's ``currency'' is shared capability, and a 2508-line file that only one person can edit is an illiquid asset. The legal analogy doesn't map perfectly, but the intellectual property point does: code that no one else can maintain is not a team asset; it is a personal portfolio piece that happens to live in a shared repository.

Three mechanisms for future refocusing:

\textbf{1. Build for named users, not for the abstract team.} Before starting any new tool, I will name at least one teammate who will use it and document their name in the commit message. The signal: at least 80\% of new files have a named user. The test I should have been running: if I built \texttt{simple\_simulator.py} for Robin to test his flight logic without a drone, would I have designed it differently? Yes --- I would have made the keyboard controls match his flight-script commands, not mine. The tool would have been less elegant and more useful.

\textbf{2. No solo refactors of files over 500 lines.} I will pair on every review of files over 500 lines. The signal: my pair partner can edit that file independently within 48 hours of the walkthrough. The test from this project: the \texttt{main.py} refactor (1411 to 754 lines, 6 extracted modules) was better code. But it was code that only I could maintain, which means the refactor made the codebase \textit{more} fragile from a team perspective, not less. I was optimising for code quality and accidentally degrading team resilience.

\textbf{3. Teaching sessions with measurable outcomes.} Two hours per week of explicit pair-programming or walkthrough, graded not by whether I explained something but by whether the teammate completes a handoff task within 48 hours without a second question. If after four weeks the signals are not firing, the mechanism must change. That is a sharper test than ``share knowledge more,'' which is the sentence that appears in every group reflection and changes nothing.

Applying the risk management framework from Graham's lectures: each of these mechanisms is a mitigation. Mechanism 1 reduces the \textit{likelihood} of building unused tools. Mechanism 2 reduces the \textit{consequence} of a team member leaving (no single-owner files). Mechanism 3 reduces both --- teaching transfers knowledge (likelihood) and creates backup capability (consequence). The residual risk: even with these mechanisms, I may still default to solo work under time pressure. That residual risk is not zero, and I acknowledge it rather than pretending the mechanisms are sufficient.

I leave this section with a self-evaluation that applies Graham's ALARP principle. Is my current approach to team contribution ``as low as reasonably practicable'' in terms of team risk? No. The risk of single-owner modules is tolerable but not negligible: if I had been ill during the last two weeks of the project, the team would have lost access to \texttt{main.py}, \texttt{simple\_simulator.py}, \texttt{passive\_watch.py}, \texttt{pi\_flight.py}, and the entire vision pipeline. That is an intolerable bus factor. The mitigations I have proposed --- named users, paired refactors, teaching sessions --- reduce the risk to tolerable. But they require me to accept slower initial progress in exchange for faster recovery from disruption, and that trade-off is the one I find hardest to make.
